\documentclass[11pt]{article}

\usepackage[margin=1in]{geometry}
\usepackage{amsmath, amssymb, bm}
\usepackage{natbib}
\usepackage{enumitem}
\usepackage{booktabs}

\newcommand{\bsy}{\bm{y}}
\newcommand{\bsx}{\bm{x}}
\newcommand{\bsb}{\bm{\beta}}
\newcommand{\bsw}{\bm{\omega}}
\newcommand{\bse}{\bm{\epsilon}}
\newcommand{\bstheta}{\bm{\theta}}

\title{Gaussian Process Change-of-Support Model: Model, Fitting, Recovery, Prediction, and Tapered-Covariance Extension}
\author{}
\date{}

\begin{document}
\maketitle

\section{Purpose and intuition}

The model is designed for the forest-inventory setting in which the field response and remotely sensed auxiliary variables are observed on different spatial supports. LiDAR-derived predictors can be summarized on a fine grid, while the field response is observed on larger fixed-area plots, and final inferential targets may be irregular stands or management units. Rather than forcing all quantities onto one arbitrary common support before modeling, we define a latent forest response surface at the fine predictor support and connect observed plots and prediction units to that surface through area-weighted averages.

Following the notation of Zhang et al., let $A_i$, $i=1,\ldots,n_a$, denote fine grid cells, let $B_l$, $l=1,\ldots,n_b$, denote observed areal units, and let $U_k$, $k=1,\ldots,n_u$, denote prediction units, such as stands, management units, or holdout polygons.

The key idea is:
$$
\text{observed plot value} = \text{area-weighted average of a fine-scale response surface over the plot}.
$$
Similarly,
$$
\text{stand value} = \text{area-weighted average of the same fine-scale response surface over the stand}.
$$
Thus, plots and stands are treated symmetrically: both are averages of the fine-scale process over their respective footprints.

\section{Fine-scale model}

Let $\bsx(A_i)$ be the $p$-dimensional predictor vector for fine pixel $A_i$, including the intercept if desired. Stack these into the $n_a \times p$ matrix
\begin{equation}
\bm{X}
=
\begin{bmatrix}
\bsx(A_1)^\top\\
\vdots\\
\bsx(A_{n_a})^\top
\end{bmatrix}.
\label{eq:X}
\end{equation}

At the fine support, the full-GP model is
\begin{equation}
y(A_i)
=
\bsx(A_i)^\top\bsb
+
\omega(A_i)
+
\epsilon(A_i),
\qquad i=1,\ldots,n_a.
\label{eq:fine_model}
\end{equation}
Here $y(A_i)$ is the fine-scale forest response density, $\bsx(A_i)^\top\bsb$ is the regression component, $\omega(A_i)$ is spatially structured residual response variation, and $\epsilon(A_i)$ is independent fine-scale noise at the finite grid-cell support.

Let
\begin{equation}
\bsw_A =
\begin{bmatrix}
\omega(A_1)\\
\vdots\\
\omega(A_{n_a})
\end{bmatrix},
\qquad
\bse_A =
\begin{bmatrix}
\epsilon(A_1)\\
\vdots\\
\epsilon(A_{n_a})
\end{bmatrix}.
\label{eq:w_e_vectors}
\end{equation}
The full-GP code uses
\begin{equation}
\bsw_A \sim N(\bm{0},\sigma^2 \bm{C}_{\phi}),
\label{eq:wA_prior}
\end{equation}
where
\begin{equation}
\left[\bm{C}_{\phi}\right]_{ij}
=
\exp\{-\phi d(s_i,s_j)\}.
\label{eq:full_exp_cor}
\end{equation}
Here $s_i$ and $s_j$ are fine-pixel centroids, and $d(s_i,s_j)$ is Euclidean distance. The parameter $\sigma^2$ controls the variance of the spatial residual process, and $\phi$ controls the decay of spatial correlation. For the exponential correlation, an approximate effective range is $3/\phi$.

The independent finite-support noise is
\begin{equation}
\bse_A \sim N(\bm{0},\tau^2\bm{I}_{n_a}).
\label{eq:eA_prior}
\end{equation}
This noise term is defined at the finite grid-cell support, not as an integral of point-level white noise.

\section{Observed areal model}

Each observed plot $B_l$ is related to the fine support through weights
\begin{equation}
h_{li}
=
\frac{|B_l \cap A_i|}{|B_l|},
\label{eq:hli}
\end{equation}
where $|\cdot|$ denotes area. Thus $h_{li}$ is the proportion of plot $B_l$'s area contributed by fine pixel $A_i$. Most $h_{li}$'s are zero because a plot overlaps only a small subset of fine pixels.

Let $\bm{H}_{BA}$ be the $n_b \times n_a$ matrix with entries $h_{li}$. Rows of $\bm{H}_{BA}$ generally sum to one if the fine grid fully covers each observed support.

The observed response on plot $B_l$ is approximated by
\begin{equation}
y(B_l)
\approx
\sum_{i=1}^{n_a} h_{li} y(A_i).
\label{eq:areal_average}
\end{equation}
Stacking observed responses into $\bsy_B$, we obtain
\begin{equation}
\bsy_B
=
\bm{H}_{BA}\bm{X}\bsb
+
\bm{H}_{BA}\bsw_A
+
\bm{H}_{BA}\bse_A.
\label{eq:yB_fine_stack}
\end{equation}

Define the areal spatial random effects
\begin{equation}
\bsw_B = \bm{H}_{BA}\bsw_A.
\label{eq:wB_def}
\end{equation}
Because $\bsw_A$ is Gaussian,
\begin{equation}
\bsw_B \sim N(\bm{0},\sigma^2\bm{C}_B(\phi)),
\label{eq:wB_prior}
\end{equation}
where
\begin{equation}
\bm{C}_B(\phi)
=
\bm{H}_{BA}\bm{C}_{\phi}\bm{H}_{BA}^\top.
\label{eq:CB_full}
\end{equation}

Similarly,
\begin{equation}
\bm{H}_{BA}\bse_A
\sim
N(\bm{0},\tau^2 \bm{D}_h),
\qquad
\bm{D}_h = \bm{H}_{BA}\bm{H}_{BA}^\top.
\label{eq:Dh}
\end{equation}
If observed supports do not overlap in terms of fine pixels, $\bm{D}_h$ is diagonal. If observed supports overlap, $\bm{D}_h$ may be non-diagonal. The code uses the general expression.

The observed-support model is therefore
\begin{equation}
\bsy_B \mid \bsb,\bsw_B,\tau^2
\sim
N\left(
\bm{H}_{BA}\bm{X}\bsb+\bsw_B,\,
\tau^2\bm{D}_h
\right),
\label{eq:yB_conditional}
\end{equation}
with
\begin{equation}
\bsw_B\mid \sigma^2,\phi
\sim
N\left(\bm{0},\sigma^2\bm{C}_B(\phi)\right).
\label{eq:wB_conditional}
\end{equation}

This formulation avoids sampling the high-dimensional $\bsw_A$ during model fitting. Instead, fitting is conducted using the lower-dimensional areal effects $\bsw_B$, as in the implementation strategy described by Zhang et al.

\section{Priors}

The code uses
\begin{align}
\bsb &\sim N(\bm{\mu}_{\beta},\bm{V}_{\beta}), \label{eq:beta_prior}\\
\tau^2 &\sim \mathrm{IG}(t_{\tau},s_{\tau}), \label{eq:tau_prior}\\
\sigma^2 &\sim \mathrm{IG}(t_{\sigma},s_{\sigma}), \label{eq:sigma_prior}\\
\phi &\sim \mathrm{Unif}(a_{\phi},b_{\phi}). \label{eq:phi_prior}
\end{align}
The inverse-gamma distribution is parameterized with shape and scale:
\begin{equation}
p(x)
=
\frac{s^t}{\Gamma(t)}
x^{-(t+1)}
\exp\{-s/x\},
\qquad x>0.
\label{eq:ig_density}
\end{equation}
Although the model is written in terms of the fine-scale nugget variance $\tau^2$, this scale can be unintuitive because the observed data see the nugget only after support averaging. The code therefore places the prior on an average observed-support nugget variance,
$$
\tau_B^2 = \tau^2 \bar{d}_h,
\qquad
\bar{d}_h =
\frac{1}{n_b}\sum_{l=1}^{n_b} D_h(l,l).
$$
Here $\tau_B^2$ is exactly the average diagonal contribution of the nugget to the observed plot-support covariance, because the observed-support nugget covariance is $\tau^2\bm{D}_h$ and its average diagonal element is $\tau^2\bar{d}_h$. The conversion back to the fine-scale variance is therefore exact given this definition:
$$
\tau^2 = \frac{\tau_B^2}{\bar{d}_h}.
$$
This reparameterization does not change the model. It only makes prior specification and MCMC sampling more interpretable. The quantity $\tau_B^2$ should not be interpreted as the nugget variance for every observed plot unless all diagonal elements $D_h(l,l)$ are equal. For plot $l$, the observed-support nugget variance remains $(\tau_B^2/\bar{d}_h)D_h(l,l)$.

Let
\begin{equation}
\bstheta = (\tau^2,\sigma^2,\phi).
\label{eq:theta_def}
\end{equation}
The sampler uses the equivalent parameterization $\bstheta_B=(\tau_B^2,\sigma^2,\phi)$ and converts each sampled $\tau_B^2$ value to $\tau^2$ before evaluating likelihoods or recovery distributions.

\section{Collapsed posterior for $\bstheta$}

The fitting code first samples $\bstheta$ from the collapsed posterior $p(\bstheta\mid \bsy_B)$, where $\bsb$ and $\bsw_B$ have been integrated out.

Given the proper Gaussian prior for $\bsb$, the marginal model is
\begin{equation}
\bsy_B \mid \bstheta
\sim
N\left(
\bm{H}_{BA}\bm{X}\bm{\mu}_{\beta},\,
\bm{V}_Y(\bstheta)
\right),
\label{eq:yB_marginal}
\end{equation}
where
\begin{equation}
\bm{V}_Y(\bstheta)
=
\sigma^2\bm{C}_B(\phi)
+
\bm{H}_{BA}\bm{X}\bm{V}_{\beta}\bm{X}^{\top}\bm{H}_{BA}^{\top}
+
\tau^2\bm{D}_h.
\label{eq:VY}
\end{equation}
Therefore,
\begin{equation}
p(\bstheta\mid \bsy_B)
\propto
N\left(
\bsy_B \mid
\bm{H}_{BA}\bm{X}\bm{\mu}_{\beta},
\bm{V}_Y(\bstheta)
\right)
p(\tau^2)p(\sigma^2)p(\phi).
\label{eq:theta_collapsed}
\end{equation}
When the code samples $\bstheta_B=(\tau_B^2,\sigma^2,\phi)$, the same covariance is evaluated using
$$
\bm{V}_Y(\bstheta_B)
=
\sigma^2\bm{C}_B(\phi)
+
\bm{H}_{BA}\bm{X}\bm{V}_{\beta}\bm{X}^{\top}\bm{H}_{BA}^{\top}
+
\frac{\tau_B^2}{\bar{d}_h}\bm{D}_h.
$$
The prior contribution is then $p(\tau_B^2)p(\sigma^2)p(\phi)$, and each proposed $\tau_B^2$ is converted to $\tau^2=\tau_B^2/\bar{d}_h$ before forming $\bm{V}_Y$.

The code evaluates this collapsed density using a Cholesky decomposition of $\bm{V}_Y(\bstheta)$, an $n_b\times n_b$ matrix. The full fine-pixel covariance $\bm{C}_{\phi}$ is constructed to form $\bm{C}_B(\phi)=\bm{H}_{BA}\bm{C}_{\phi}\bm{H}_{BA}^{\top}$, but $\bm{C}_{\phi}$ is not inverted during fitting.

\section{Metropolis update for $\bstheta$}

The code updates $\bstheta_B=(\tau_B^2,\sigma^2,\phi)$ as a block on an unconstrained scale. The transformations are
\begin{equation}
z_1 = \log \tau_B^2,
\qquad
z_2 = \log \sigma^2,
\qquad
z_3
=
\log\left(\frac{\phi-a_{\phi}}{b_{\phi}-\phi}\right).
\label{eq:z_transform}
\end{equation}
The inverse transformation is
\begin{equation}
\tau_B^2 = \exp(z_1),
\qquad
\tau^2 = \tau_B^2/\bar{d}_h,
\qquad
\sigma^2 = \exp(z_2),
\qquad
\phi
=
a_{\phi}
+
(b_{\phi}-a_{\phi})
\frac{\exp(z_3)}{1+\exp(z_3)}.
\label{eq:z_inverse}
\end{equation}

A random-walk proposal is generated as
\begin{equation}
\bm{z}^{*}
=
\bm{z}
+
\bm{\delta},
\qquad
\bm{\delta}\sim N(\bm{0},\bm{\Sigma}_q),
\label{eq:rw_proposal}
\end{equation}
where the current implementation uses diagonal $\bm{\Sigma}_q$. The proposal is accepted with probability
\begin{equation}
\alpha
=
\min\left\{
1,\,
\frac{
p(\bstheta_B(\bm{z}^{*})\mid \bsy_B)
\left|\frac{\partial \bstheta_B}{\partial \bm{z}^{*}}\right|
}{
p(\bstheta_B(\bm{z})\mid \bsy_B)
\left|\frac{\partial \bstheta_B}{\partial \bm{z}}\right|
}
\right\}.
\label{eq:mh_alpha}
\end{equation}

The log-Jacobian is
\begin{equation}
\log\left|\frac{\partial \bstheta_B}{\partial \bm{z}}\right|
=
\log\tau_B^2
+
\log\sigma^2
+
\log(b_{\phi}-a_{\phi})
+
\log p_{\phi}
+
\log(1-p_{\phi}),
\label{eq:jacobian}
\end{equation}
where
\begin{equation}
p_{\phi}
=
\frac{\phi-a_{\phi}}{b_{\phi}-a_{\phi}}.
\label{eq:p_phi}
\end{equation}
This Jacobian is included in the code.

\section{Recovery of $\bsb$ and $\bsw_B$}

After retaining posterior samples $\{ \tau_B^{2(g)},\sigma^{2(g)},\phi^{(g)} \}_{g=1}^{G}$, the code first converts each sampled observed-support nugget variance to the corresponding fine-scale nugget variance,
$$
\tau^{2(g)} = \frac{\tau_B^{2(g)}}{\bar{d}_h}.
$$
The code then recovers one draw of $(\bsb,\bsw_B)$ for each retained $\bstheta^{(g)}=(\tau^{2(g)},\sigma^{2(g)},\phi^{(g)})$.

Let
\begin{equation}
\bm{X}_B = \bm{H}_{BA}\bm{X},
\qquad
\bm{Z} = [\bm{X}_B,\bm{I}_{n_b}],
\label{eq:Z_def}
\end{equation}
and define
\begin{equation}
\bm{\alpha}
=
\begin{bmatrix}
\bsb\\
\bsw_B
\end{bmatrix}.
\label{eq:alpha_def}
\end{equation}
Then
\begin{equation}
\bsy_B \mid \bm{\alpha},\tau^2
\sim
N(\bm{Z}\bm{\alpha},\tau^2\bm{D}_h).
\label{eq:yB_alpha}
\end{equation}

The prior for $\bm{\alpha}$ is block Gaussian:
\begin{equation}
\bm{\alpha}
\sim
N\left(
\begin{bmatrix}
\bm{\mu}_{\beta}\\
\bm{0}
\end{bmatrix},
\begin{bmatrix}
\bm{V}_{\beta} & \bm{0}\\
\bm{0} & \sigma^2\bm{C}_B(\phi)
\end{bmatrix}
\right).
\label{eq:alpha_prior}
\end{equation}

Thus,
\begin{equation}
\bm{\alpha}\mid \bstheta,\bsy_B
\sim
N(\bm{M}\bm{m},\bm{M}),
\label{eq:alpha_cond}
\end{equation}
where
\begin{equation}
\bm{M}^{-1}
=
\bm{Z}^{\top}(\tau^2\bm{D}_h)^{-1}\bm{Z}
+
\begin{bmatrix}
\bm{V}_{\beta}^{-1} & \bm{0}\\
\bm{0} & \sigma^{-2}\bm{C}_B(\phi)^{-1}
\end{bmatrix},
\label{eq:M_inv}
\end{equation}
and
\begin{equation}
\bm{m}
=
\bm{Z}^{\top}(\tau^2\bm{D}_h)^{-1}\bsy_B
+
\begin{bmatrix}
\bm{V}_{\beta}^{-1}\bm{\mu}_{\beta}\\
\bm{0}
\end{bmatrix}.
\label{eq:m_vec}
\end{equation}
The code samples this Gaussian full conditional using a Cholesky factorization of the precision matrix $\bm{M}^{-1}$, whose dimension is $(p+n_b)\times(p+n_b)$.

\section{Fine-scale recovery: $\bsw_A$ and $\bm{\eta}_A$}

For mapping and stand-level estimation, the desired target is usually the latent mean surface
\begin{equation}
\bm{\eta}_A
=
\bm{X}\bsb+\bsw_A.
\label{eq:eta_A}
\end{equation}
This excludes the fine-scale nugget $\bse_A$. It represents the model-based underlying forest response density at the fine support.

The fitting algorithm samples $\bsw_B=\bm{H}_{BA}\bsw_A$, not $\bsw_A$. Under the full GP, the conditional distribution of $\bsw_A$ given $\bsw_B$, $\sigma^2$, and $\phi$ is Gaussian. Let
\begin{equation}
\bm{C}_A = \bm{C}_{\phi},
\qquad
\bm{C}_B = \bm{H}_{BA}\bm{C}_A\bm{H}_{BA}^{\top},
\qquad
\bm{C}_{AB} = \bm{C}_A\bm{H}_{BA}^{\top}.
\label{eq:CA_CB_CAB}
\end{equation}
Then
\begin{equation}
\bsw_A \mid \bsw_B,\sigma^2,\phi
\sim
N\left(
\bm{C}_{AB}\bm{C}_B^{-1}\bsw_B,\,
\sigma^2\left[
\bm{C}_A
-
\bm{C}_{AB}\bm{C}_B^{-1}\bm{C}_{AB}^{\top}
\right]
\right).
\label{eq:wA_cond}
\end{equation}

For each retained posterior draw $g$, the code can either sample
\begin{equation}
\bsw_A^{(g)} \sim
p(\bsw_A\mid \bsw_B^{(g)},\sigma^{2(g)},\phi^{(g)})
\label{eq:wA_sample}
\end{equation}
or use the conditional mean. The latter is faster and useful for exploratory mapping, but it understates uncertainty because it omits conditional uncertainty in $\bsw_A$.

Given $\bsw_A^{(g)}$, the fine-scale latent mean draw is
\begin{equation}
\bm{\eta}_A^{(g)}
=
\bm{X}\bsb^{(g)}+\bsw_A^{(g)}.
\label{eq:eta_A_draw}
\end{equation}

If a noisy fine-scale posterior predictive draw is desired, one may additionally sample
\begin{equation}
\bm{y}_A^{(g)}
=
\bm{\eta}_A^{(g)}
+
\bse_A^{(g)},
\qquad
\bse_A^{(g)}\sim N(\bm{0},\tau^{2(g)}\bm{I}_{n_a}).
\label{eq:yA_pred}
\end{equation}
However, for mapping the underlying forest condition and for estimating stand means/totals, $\bm{\eta}_A$ is usually the target.

\section{Prediction for areal units}

Let $U_k$, $k=1,\ldots,n_u$, denote prediction units, such as stands. Define
\begin{equation}
g_{ki}
=
\frac{|U_k\cap A_i|}{|U_k|},
\label{eq:gki}
\end{equation}
and let $\bm{H}_{UA}$ be the $n_u\times n_a$ matrix with entries $g_{ki}$. Then the latent prediction-unit mean is
\begin{equation}
\bm{\eta}_U
=
\bm{H}_{UA}\bm{\eta}_A.
\label{eq:eta_U}
\end{equation}
For posterior draw $g$,
\begin{equation}
\bm{\eta}_U^{(g)}
=
\bm{H}_{UA}\bm{\eta}_A^{(g)}.
\label{eq:eta_U_draw}
\end{equation}

The expression in Equation~\eqref{eq:eta_U_draw} is useful conceptually because it makes the prediction target explicit: a prediction-unit mean is an area-weighted average of the fine-scale latent mean surface. For computation, however, it is not always necessary to recover $\bm{\eta}_A$ over every fine pixel. Prediction can also be done directly at the areal support.

Define the prediction-unit spatial random effects
\begin{equation}
\bsw_U = \bm{H}_{UA}\bsw_A.
\label{eq:wU_def}
\end{equation}
Let
\begin{equation}
\bm{C}_U
=
\bm{H}_{UA}\bm{C}_A\bm{H}_{UA}^{\top},
\qquad
\bm{C}_{UB}
=
\bm{H}_{UA}\bm{C}_A\bm{H}_{BA}^{\top},
\label{eq:CU_CUB}
\end{equation}
where $\bm{C}_A=\bm{C}_{\phi}$, and let $\bm{C}_{BU}=\bm{C}_{UB}^{\top}$. Then the joint Gaussian model implies
\begin{equation}
\bsw_U \mid \bsw_B,\sigma^2,\phi
\sim
N\left(
\bm{C}_{UB}\bm{C}_B^{-1}\bsw_B,\,
\sigma^2\left[
\bm{C}_U
-
\bm{C}_{UB}\bm{C}_B^{-1}\bm{C}_{BU}
\right]
\right).
\label{eq:wU_cond}
\end{equation}
Thus, for posterior draw $g$, one may draw $\bsw_U^{(g)}$ from Equation~\eqref{eq:wU_cond} and compute
\begin{equation}
\bm{\eta}_U^{(g)}
=
\bm{H}_{UA}\bm{X}\bsb^{(g)} + \bsw_U^{(g)}.
\label{eq:eta_U_direct_draw}
\end{equation}
If only the conditional mean spatial adjustment is desired, then
\begin{equation}
E\{\bm{\eta}_U^{(g)} \mid \bsw_B^{(g)},\bstheta^{(g)}\}
=
\bm{H}_{UA}\bm{X}\bsb^{(g)}
+
\bm{C}_{UB}^{(g)}\{\bm{C}_B^{(g)}\}^{-1}\bsw_B^{(g)}.
\label{eq:eta_U_direct_mean}
\end{equation}
This direct-to-stand form is algebraically equivalent to aggregating the fine-scale process when the same covariance model is used. Its advantage is computational: for stand-level inference, it can avoid storing or sampling $\bm{\eta}_A$ over a very large fine grid. The fine grid is still used to construct $\bm{H}_{UA}\bm{X}$, $\bm{C}_{UB}$, and $\bm{C}_U$, but posterior samples can be retained at the stand support rather than at every fine pixel.

If $U_k$ has area $|U_k|_{\mathrm{ha}}$ in hectares and $\eta_U$ is in response units per hectare, then the total response in $U_k$ is
\begin{equation}
T_k^{(g)}
=
|U_k|_{\mathrm{ha}}\eta_U^{(g)}(k).
\label{eq:total_U}
\end{equation}
Posterior summaries of $\{\eta_U^{(g)}(k):g=1,\ldots,G\}$ give the posterior distribution for stand mean response. Posterior summaries of $\{T_k^{(g)}:g=1,\ldots,G\}$ give the posterior distribution for stand total response.

\section{Posterior predictive distribution for a future areal observation}

The latent stand mean $\eta_U$ is not the same as a noisy future observation. If the target is a posterior predictive distribution for a new areal observation on support $U$, then the nugget contribution should be included:
\begin{equation}
\bm{y}_U
=
\bm{\eta}_U+\bm{e}_U,
\qquad
\bm{e}_U\sim N(\bm{0},\tau^2\bm{D}_U),
\label{eq:yU_pred}
\end{equation}
where
\begin{equation}
\bm{D}_U
=
\bm{H}_{UA}\bm{H}_{UA}^{\top}.
\label{eq:DU}
\end{equation}
Thus, for draw $g$,
\begin{equation}
\bm{y}_U^{(g)}
\sim
N\left(
\bm{\eta}_U^{(g)},
\tau^{2(g)}\bm{D}_U
\right).
\label{eq:yU_draw}
\end{equation}

The distinction is important. $\eta_U$ is the stand's underlying mean response and is usually the target for inventory-style stand estimates. $y_U$ is a noisy posterior predictive realization of a future areal measurement.

\section{Algorithm summary}

During fitting:
\begin{enumerate}[leftmargin=2em]
\item Given current transformed parameters $\bm{z}$, propose $\bm{z}^{*}$ using the block random-walk proposal in Equation~\eqref{eq:rw_proposal}.
\item Transform $\bm{z}^{*}$ to $\bstheta_B^{*}=(\tau_B^{2*},\sigma^{2*},\phi^{*})$ using Equation~\eqref{eq:z_inverse}, then compute $\tau^{2*}=\tau_B^{2*}/\bar{d}_h$ for evaluating the likelihood.
\item Construct $\bm{C}_{\phi^{*}}$ using Equation~\eqref{eq:full_exp_cor} and $\bm{C}_B(\phi^{*})$ using Equation~\eqref{eq:CB_full}.
\item Evaluate the collapsed log posterior in Equation~\eqref{eq:theta_collapsed}, including the transformation Jacobian in Equation~\eqref{eq:jacobian}.
\item Accept or reject $\bstheta_B^{*}$ using the Metropolis ratio in Equation~\eqref{eq:mh_alpha}.
\end{enumerate}

After retaining $G$ samples of $\bstheta_B$:
\begin{enumerate}[leftmargin=2em]
\item Convert each retained $\tau_B^{2(g)}$ to $\tau^{2(g)}=\tau_B^{2(g)}/\bar{d}_h$.
\item For each retained $\bstheta^{(g)}=(\tau^{2(g)},\sigma^{2(g)},\phi^{(g)})$, sample $(\bsb^{(g)},\bsw_B^{(g)})$ from Equation~\eqref{eq:alpha_cond}.
\item Optionally recover or sample $\bsw_A^{(g)}$ from Equation~\eqref{eq:wA_cond}.
\item Compute $\bm{\eta}_A^{(g)}$ using Equation~\eqref{eq:eta_A_draw}.
\item For each prediction-unit matrix $\bm{H}_{UA}$, compute $\bm{\eta}_U^{(g)}$ either by aggregating the recovered fine-scale surface using Equation~\eqref{eq:eta_U_draw}, or directly at the prediction-unit support using Equation~\eqref{eq:eta_U_direct_draw} or Equation~\eqref{eq:eta_U_direct_mean}.
\item If totals are required, compute $T_k^{(g)}$ using Equation~\eqref{eq:total_U}.
\item If a noisy posterior predictive areal observation is required, sample $\bm{y}_U^{(g)}$ using Equation~\eqref{eq:yU_draw}.
\end{enumerate}

\subsection*{Tapered version of the algorithm}

For the tapered implementation, the algorithm is unchanged except for the covariance construction in Step 3. The taper distance $\gamma$ is fixed before running the sampler. Step 3 is replaced by: construct the tapered covariance entries $\bm{C}^{\mathrm{tap}}_{\phi^*,\gamma}$ using Equation~\eqref{eq:tapered_cor}, or preferably compute $\bm{C}^{\mathrm{tap}}_B(\phi^*,\gamma)$ directly using Equation~\eqref{eq:CB_taper_double_sum}, without forming the full fine-scale matrix. All subsequent likelihood and recovery calculations use $\bm{C}^{\mathrm{tap}}_B(\phi^*,\gamma)$ in place of $\bm{C}_B(\phi^*)$. The Metropolis proposal remains a block update for $\bstheta_B=(\tau_B^2,\sigma^2,\phi)$; $\gamma$ is not proposed or sampled.

\section{Implementation notes for the current full-GP code}

The current code is intentionally a full-GP reference implementation. It is useful for checking model algebra, simulation behavior, posterior recovery, and prediction mechanics. It is not intended to scale directly to very large fine grids.

During fitting, the largest matrix decompositions are the $n_b\times n_b$ Cholesky decomposition of $\bm{V}_Y(\bstheta)$ in Equation~\eqref{eq:VY} and the $(p+n_b)\times(p+n_b)$ Cholesky decomposition used to sample Equation~\eqref{eq:alpha_cond}. The code does not invert or decompose the $n_a\times n_a$ matrix $\bm{C}_{\phi}$ during fitting. The computational bottleneck in the full-GP fitting code is constructing the dense matrix $\bm{C}_{\phi}$ in Equation~\eqref{eq:full_exp_cor} and forming $\bm{C}_B(\phi)$ in Equation~\eqref{eq:CB_full} for each proposed $\phi$.

During fine-scale recovery, sampling $\bsw_A\mid \bsw_B$ from Equation~\eqref{eq:wA_cond} requires operations involving the full $n_a\times n_a$ conditional covariance. This is feasible only for small development examples. For large fine grids, fine-scale recovery and prediction should be implemented using the same scalable approximation used for fitting.

\section{Tapered covariance extension}

The full-GP implementation uses the dense fine-scale correlation in Equation~\eqref{eq:full_exp_cor}. This is conceptually simple, but it is not scalable when $n_a$ is large. A tapered covariance provides a direct way to retain the same change-of-support model while avoiding construction of the full $n_a\times n_a$ covariance matrix.

In the tapered version considered here, the taper distance $\gamma$ is treated as a fixed hyperparameter, chosen before running the sampler. It is not included in $\bstheta=(\tau^2,\sigma^2,\phi)$ and is not updated in the MCMC algorithm. Thus, the posterior sampling steps remain the same as in the full-GP implementation, except that all covariance terms involving $\bm{C}_{\phi}$ are replaced by their tapered analogs involving $\bm{C}^{\mathrm{tap}}_{\phi,\gamma}$.

\subsection{Tapered fine-scale covariance}

Let $T_{\gamma}(d)$ be a compactly supported taper function with taper distance $\gamma$, such that
\begin{equation}
T_{\gamma}(d)=0
\qquad \text{for } d\ge \gamma.
\label{eq:taper_zero}
\end{equation}
A common choice is the Wendland taper
\begin{equation}
T_{\gamma}(d)
=
\left(1-\frac{d}{\gamma}\right)^4
\left(1+4\frac{d}{\gamma}\right)
\mathbb{I}(d<\gamma).
\label{eq:wendland}
\end{equation}
The tapered version replaces the full exponential correlation in Equation~\eqref{eq:full_exp_cor} with
\begin{equation}
\left[\bm{C}^{\mathrm{tap}}_{\phi,\gamma}\right]_{ij}
=
\exp\{-\phi d(s_i,s_j)\}
T_{\gamma}\{d(s_i,s_j)\}.
\label{eq:tapered_cor}
\end{equation}

Intuitively, the taper says that fine pixels farther apart than $\gamma$ have zero residual spatial covariance. Nearby pixels still have covariance governed by the exponential correlation, but that covariance is smoothly pushed to zero as distance approaches $\gamma$. If the LiDAR predictors explain most broad-scale structure and the residual process has short range, this is a reasonable approximation and can be much more efficient.

\subsection{Where the tapered covariance is swapped into the model}

The model structure in Equation~\eqref{eq:fine_model} does not change. The observation weights in Equations~\eqref{eq:hli}--\eqref{eq:areal_average} do not change. The only change is the covariance used for the spatial random effect.

Specifically, replace the full-GP prior in Equation~\eqref{eq:wA_prior} with
\begin{equation}
\bsw_A
\sim
N(\bm{0},\sigma^2\bm{C}^{\mathrm{tap}}_{\phi,\gamma}).
\label{eq:wA_taper_prior}
\end{equation}
Then Equation~\eqref{eq:CB_full} becomes
\begin{equation}
\bm{C}^{\mathrm{tap}}_B(\phi,\gamma)
=
\bm{H}_{BA}
\bm{C}^{\mathrm{tap}}_{\phi,\gamma}
\bm{H}_{BA}^{\top}.
\label{eq:CB_taper}
\end{equation}
The observed-support spatial prior in Equation~\eqref{eq:wB_conditional} becomes
\begin{equation}
\bsw_B\mid \sigma^2,\phi,\gamma
\sim
N\left(
\bm{0},
\sigma^2\bm{C}^{\mathrm{tap}}_B(\phi,\gamma)
\right).
\label{eq:wB_taper_cond}
\end{equation}
The conditioning on $\gamma$ emphasizes that the covariance depends on the chosen taper distance; $\gamma$ is fixed by the analyst rather than sampled as part of $\bstheta$.
The marginal covariance in Equation~\eqref{eq:VY} becomes
\begin{equation}
\bm{V}^{\mathrm{tap}}_Y(\bstheta)
=
\sigma^2\bm{C}^{\mathrm{tap}}_B(\phi,\gamma)
+
\bm{H}_{BA}\bm{X}\bm{V}_{\beta}\bm{X}^{\top}\bm{H}_{BA}^{\top}
+
\tau^2\bm{D}_h.
\label{eq:VY_taper}
\end{equation}
The collapsed posterior in Equation~\eqref{eq:theta_collapsed} is then evaluated using $\bm{V}^{\mathrm{tap}}_Y(\bstheta)$ in place of $\bm{V}_Y(\bstheta)$.

The recovery of $(\bsb,\bsw_B)$ in Equations~\eqref{eq:alpha_cond}--\eqref{eq:m_vec} is unchanged except that $\bm{C}_B(\phi)$ is replaced by $\bm{C}^{\mathrm{tap}}_B(\phi,\gamma)$.

Fine-scale recovery in Equation~\eqref{eq:wA_cond} is similarly modified by replacing
$$
\bm{C}_A=\bm{C}_{\phi}
$$
with
$$
\bm{C}^{\mathrm{tap}}_A=\bm{C}^{\mathrm{tap}}_{\phi,\gamma},
$$
and replacing $\bm{C}_B$ and $\bm{C}_{AB}$ by their tapered analogs:
\begin{equation}
\bm{C}^{\mathrm{tap}}_B
=
\bm{H}_{BA}\bm{C}^{\mathrm{tap}}_A\bm{H}_{BA}^{\top},
\qquad
\bm{C}^{\mathrm{tap}}_{AB}
=
\bm{C}^{\mathrm{tap}}_A\bm{H}_{BA}^{\top}.
\label{eq:taper_prediction_covs}
\end{equation}
Equivalently, the fine-pixel to observed-support cross-covariance has entries
\begin{equation}
C^{\mathrm{tap}}_{AB}(i,l)
=
\sum_{j=1}^{n_a}
h_{lj}
\exp\{-\phi d(s_i,s_j)\}
T_{\gamma}\{d(s_i,s_j)\}.
\label{eq:CAB_taper_double_sum}
\end{equation}
Only pixels $A_j$ that overlap $B_l$ and are within distance $\gamma$ of $A_i$ contribute to this sum.

\subsection{How the computation changes}

The important computational point is that the tapered model need not explicitly form the dense matrix $\bm{C}^{\mathrm{tap}}_{\phi,\gamma}$. Instead, each areal covariance entry can be computed directly as
\begin{equation}
C^{\mathrm{tap}}_B(l,k)
=
\sum_{i=1}^{n_a}
\sum_{j=1}^{n_a}
h_{li}h_{kj}
\exp\{-\phi d(s_i,s_j)\}
T_{\gamma}\{d(s_i,s_j)\}.
\label{eq:CB_taper_double_sum}
\end{equation}
Because $h_{li}$ and $h_{kj}$ are zero for most fine pixels, and because $T_{\gamma}(d)=0$ for $d\ge\gamma$, Equation~\eqref{eq:CB_taper_double_sum} only requires local fine-pixel pairs that both contribute to the relevant areal supports and are within the taper distance.

This is where the taper helps. The full-GP code computes $\bm{C}_{\phi}$ for all fine-pixel pairs, even though the sampler only needs the $n_b\times n_b$ induced covariance $\bm{C}_B(\phi)$. The intended tapered implementation computes $\bm{C}^{\mathrm{tap}}_B(\phi,\gamma)$ directly from local pairwise contributions. Thus, the matrix decompositions in the sampler remain $n_b\times n_b$ and $(p+n_b)\times(p+n_b)$, but the expensive dense $n_a\times n_a$ covariance construction is avoided.

The computational savings described here require computing $\bm{C}^{\mathrm{tap}}_B(\phi,\gamma)$ directly from local pairwise contributions, as in Equation~\eqref{eq:CB_taper_double_sum}. If one instead forms the full $n_a\times n_a$ tapered matrix first, even as a sparse matrix, the savings depend on the sparsity induced by $\gamma$ and may still be unnecessarily expensive.

For prediction to units $U$, the same idea applies. The cross-covariance between fine-scale or prediction-unit quantities and observed plot effects is computed using local tapered covariance contributions rather than by multiplying full dense covariance matrices. For example, the covariance between a prediction unit $U_k$ and an observed plot $B_l$ is
\begin{equation}
C^{\mathrm{tap}}_{UB}(k,l)
=
\sum_{i=1}^{n_a}
\sum_{j=1}^{n_a}
g_{ki}h_{lj}
\exp\{-\phi d(s_i,s_j)\}
T_{\gamma}\{d(s_i,s_j)\}.
\label{eq:CUB_taper}
\end{equation}
This replaces the dense-matrix expression $\bm{H}_{UA}\bm{C}_{\phi}\bm{H}_{BA}^{\top}$. Again, only local pairs within distance $\gamma$ contribute. The direct stand-level prediction equations in Equations~\eqref{eq:wU_cond}--\eqref{eq:eta_U_direct_mean} are then used with $\bm{C}_{UB}^{\mathrm{tap}}$, $\bm{C}_{B}^{\mathrm{tap}}$, and, if full joint stand-level draws are required, $\bm{C}_{U}^{\mathrm{tap}}=\bm{H}_{UA}\bm{C}^{\mathrm{tap}}_{\phi,\gamma}\bm{H}_{UA}^{\top}$.

\subsection{Interpretation of the taper distance}

The taper distance $\gamma$ is not the same as the spatial range parameter $1/\phi$ or the effective range $3/\phi$. The parameter $\phi$ controls the correlation decay inside the support of the taper. The taper distance $\gamma$ determines where the covariance is forced exactly to zero. In practice, $\gamma$ should be chosen large enough that correlations beyond $\gamma$ are scientifically negligible or adequately captured by the fixed effects and other model terms.

If the residual spatial dependence after accounting for LiDAR predictors is short range, a relatively small $\gamma$ can be appropriate and computationally efficient. If the residual process has long-range structure, too small a taper distance can bias spatial dependence downward and understate uncertainty. Thus, $\gamma$ should be treated as a modeling and computational choice, and sensitivity to plausible taper distances should be examined.

\subsection{What remains unchanged}

The taper does not change the support logic. The weights $h_{li}$ in Equation~\eqref{eq:hli} and $g_{ki}$ in Equation~\eqref{eq:gki} are unchanged. The observed response is still an areal average as in Equation~\eqref{eq:areal_average}. The prediction target $\bm{\eta}_U=\bm{H}_{UA}\bm{\eta}_A$ in Equation~\eqref{eq:eta_U} is unchanged. The Metropolis update, conditional recovery of $(\bsb,\bsw_B)$, and posterior summaries all retain the same structure; only the covariance calculations in Equations~\eqref{eq:CB_full}, \eqref{eq:VY}, \eqref{eq:wA_cond}, and the prediction cross-covariances are replaced by their tapered counterparts.

\end{document}